\documentclass[11pt,a4paper]{article}

\usepackage[margin=2.2cm,top=2cm,bottom=2cm]{geometry}
\usepackage{booktabs}
\usepackage{array}
\usepackage{tabularx}
\usepackage{longtable}
\usepackage{xcolor}
\usepackage{listings}
\usepackage{upquote}
\usepackage{textcomp}
\usepackage{enumitem}
\usepackage{titlesec}
\usepackage{mdframed}
\usepackage{fancyhdr}
\usepackage[colorlinks=true,linkcolor=blue!45!black,urlcolor=blue!60!black]{hyperref}

\setlength{\parindent}{0pt}
\setlength{\parskip}{5pt}
\titlespacing*{\section}{0pt}{14pt}{6pt}
\titlespacing*{\subsection}{0pt}{10pt}{4pt}
\titlespacing*{\subsubsection}{0pt}{8pt}{3pt}
\setlist[itemize]{leftmargin=1.4em,itemsep=2pt,topsep=3pt}
\setlist[enumerate]{leftmargin=1.8em,itemsep=2pt,topsep=3pt}
\renewcommand{\arraystretch}{1.15}

\newcommand{\code}[1]{\texttt{\small #1}}
\newcolumntype{L}[1]{>{\raggedright\arraybackslash}p{#1}}

\newmdenv[linewidth=1pt,linecolor=blue!40!black,backgroundcolor=blue!4,
  innertopmargin=7pt,innerbottommargin=7pt,innerleftmargin=9pt,
  innerrightmargin=9pt,skipabove=8pt,skipbelow=8pt]{taskbox}
\newmdenv[linewidth=1pt,linecolor=orange!60!black,backgroundcolor=orange!6,
  innertopmargin=7pt,innerbottommargin=7pt,innerleftmargin=9pt,
  innerrightmargin=9pt,skipabove=8pt,skipbelow=8pt]{warnbox}
\newmdenv[linewidth=0.8pt,linecolor=black!35,backgroundcolor=black!3,
  innertopmargin=6pt,innerbottommargin=6pt,innerleftmargin=8pt,
  innerrightmargin=8pt,skipabove=6pt,skipbelow=6pt]{notebox}

\lstdefinestyle{sh}{basicstyle=\ttfamily\small, breaklines=true,
  columns=fullflexible, backgroundcolor=\color{black!5}, frame=single,
  framesep=5pt, rulecolor=\color{black!30}, xleftmargin=2pt,
  showstringspaces=false, aboveskip=6pt, belowskip=6pt, upquote=true}
\lstdefinestyle{sql}{basicstyle=\ttfamily\small, language=SQL, breaklines=true,
  columns=fullflexible, keywordstyle=\color{blue!60!black}\bfseries,
  commentstyle=\color{green!35!black}\itshape, stringstyle=\color{red!55!black},
  backgroundcolor=\color{black!5}, frame=single, framesep=5pt,
  rulecolor=\color{black!30}, xleftmargin=2pt, showstringspaces=false,
  aboveskip=6pt, belowskip=6pt}

\pagestyle{fancy}\fancyhf{}
\rhead{\small Benchmark Evaluation}
\lhead{\small TPC-DI \& MIMIC-IV$\rightarrow$OMOP}
\cfoot{\thepage}

\title{\vspace{-1.5cm}\textbf{Benchmark Evaluation: TPC-DI and MIMIC-IV $\rightarrow$ OMOP}\\[4pt]
\large Scope, known findings, and work plan}
\author{}
\date{}

\begin{document}
\maketitle
\vspace{-1.2cm}



% ======================================================================
\section{Background and scope}

\subsection{The research problem}

Companies frequently have to move data from one database to another --- for example, when
two companies merge, or when an old system is replaced. The old database and the new one
usually store the same information in \emph{different shapes}. Someone has to write a
program that reads the old data and writes it into the new shape. That program is called a
\textbf{transformation} (or \textbf{ETL}, see glossary).

Our research concerns generating such programs automatically and verifying their
correctness. That requires examples to evaluate against; this document covers sourcing and
preparing them.

\subsection{The four things we need}

For each example we need all four of these. If any one is missing, the example is not
usable.

\begin{center}
\begin{tabular}{@{}clL{9.2cm}@{}}
\toprule
\textbf{Symbol} & \textbf{Name} & \textbf{What it is} \\
\midrule
$S$ & Source schema & The \emph{structure} of the old database --- which tables exist, which columns they have, what type each column is. Written as SQL \code{CREATE TABLE} statements. \\
$D$ & Source data & The actual rows sitting in the old database. \\
$T$ & Target schema & The structure of the new database. Also \code{CREATE TABLE} statements. \\
--- & Transformation & The program (usually SQL) that reads $D$ and fills the new database. This is the ``correct answer'' we compare against. \\
\bottomrule
\end{tabular}
\end{center}

Two candidate sources are under evaluation:

\begin{enumerate}
\item \textbf{TPC-DI} --- a financial-trading benchmark. Section~\ref{sec:tpcdi}.
\item \textbf{MIMIC-IV $\rightarrow$ OMOP} --- a hospital-data conversion. Section~\ref{sec:mimic}.
\end{enumerate}

\subsection{Terminology}

\begin{center}
\small
\begin{longtable}{@{}L{3.3cm}L{12.4cm}@{}}
\toprule
\textbf{Term} & \textbf{Meaning} \\
\midrule
\endhead
\textbf{DDL} & ``Data Definition Language''. The SQL commands that create structure, i.e.\ \code{CREATE TABLE}. As opposed to commands that read or write data. \\
\textbf{Surrogate key} & An identifier \emph{invented} by the transformation, which does not exist in the source. Usually just 1, 2, 3, \ldots. Important: two correct transformations may number them differently, and \emph{both are right}. \\
\textbf{ETL} & ``Extract, Transform, Load''. The general name for a program that reads data from one place, changes its shape, and writes it somewhere else. \\
\textbf{Star schema} & A common warehouse design. A few large ``fact'' tables recording events (trades, sales), surrounded by smaller ``dimension'' tables describing things (customers, dates). Named because a diagram looks like a star. \\
\textbf{SCD Type 2} & ``Slowly Changing Dimension, Type 2''. A way of keeping history. When a customer moves house, you do \emph{not} overwrite the old address. You mark the old row as finished (\code{EndDate}) and add a new row marked current (\code{IsCurrent = true}). One real customer becomes several rows. \\
\textbf{DuckDB} & A database that lives in one file on your laptop. No server to install and run --- one \code{pip install} and you are done. We use it because it is the least setup. \\
\textbf{dbt} & A tool where each output table is defined by one SQL file. Useful to us because it means the transformation for each table is a separate, readable file rather than one huge script. \\
\textbf{Scale factor (SF)} & A dial on a data generator saying how much data to make. SF=10 for TPC-DI is roughly 1\,GB. Always start at the smallest. \\
\bottomrule
\end{longtable}
\end{center}

\subsection{Environment}

The toolchain assumes Linux. On Windows, use WSL.

\begin{lstlisting}[style=sh]
# 1. Open PowerShell AS ADMINISTRATOR, run this, then restart the computer:
wsl --install -d Ubuntu-22.04

# 2. After restart, open "Ubuntu" from the Start menu and set a username/password.

# 3. Inside Ubuntu, run:
sudo apt update
sudo apt install -y git python3 python3-venv python3-pip unzip

# 4. Verify:
git --version
python3 --version
\end{lstlisting}

VS Code with the WSL extension (\url{https://code.visualstudio.com/}) is convenient for
reading the SQL files.

\begin{notebox}
All work goes under \code{\textasciitilde/benchmark/}:
\begin{lstlisting}[style=sh]
mkdir -p ~/benchmark
cd ~/benchmark
\end{lstlisting}
\end{notebox}

\subsection{Record of work}

Keep \code{\textasciitilde/benchmark/FINDINGS.md} updated per task: what was run, what
happened, and the verdict. Paste error messages verbatim. This file is the primary
deliverable.

% ======================================================================
\clearpage
\section{Benchmark 1: TPC-DI}
\label{sec:tpcdi}

\subsection{What TPC-DI is}

TPC is an organisation that publishes standard tests for database systems. TPC-DI is their
test for data integration. It describes an imaginary stockbroking company:

\begin{itemize}
\item \textbf{The source} is a messy pile of files --- customer records, trade records,
  financial statements --- in several different formats. This is our $D$.
\item \textbf{The target} is a clean data warehouse with 20 tables. This is our $T$.
\item \textbf{The rules} for getting from one to the other are written in a
  specification document, in English prose. Note: \emph{prose, not code}. TPC gives the
  requirements; other people have written the actual programs.
\end{itemize}

\subsection{What we already know (verified)}

\begin{center}
\begin{tabular}{@{}L{4.2cm}L{11.3cm}@{}}
\toprule
\textbf{Finding} & \textbf{Detail} \\
\midrule
$T$ is ready & The target schema has been extracted, converted to portable SQL, and
tested. It loads cleanly on DuckDB: \textbf{20 tables, 30 foreign keys, 241 columns}.
File: \code{tpcdi\_target\_schema.sql} (supplied with this document). \\
\addlinespace[2pt]
The work is non-trivial & Of the 241 target columns, \textbf{97 (40\%) need more than a
straight copy} from the source. Breakdown: 30 need surrogate-key lookups, 15 are SCD~Type~2
history bookkeeping, 11 are derived calendar labels, 11 need lookup joins, 4 need window
functions over history, and the rest are conditionals and arithmetic. \\
\addlinespace[2pt]
Five tables keep history & \code{DimCustomer}, \code{DimCompany}, \code{DimSecurity},
\code{DimBroker}, \code{DimAccount} all use SCD~Type~2. \\
\addlinespace[2pt]
Four tables are trivial & \code{Industry}, \code{StatusType}, \code{TaxRate},
\code{TradeType} are pure copies. Useful as an easy baseline. \\
\bottomrule
\end{tabular}
\end{center}

\subsection{What is missing}

\begin{center}
\begin{tabular}{@{}cL{5.2cm}L{8.0cm}@{}}
\toprule
& \textbf{Missing piece} & \textbf{Why} \\
\midrule
$S$ & The source schema & TPC-DI's source is \emph{files}, not a database. Somebody has to write \code{CREATE TABLE} statements for each file. \\
$D$ & The source data & Must be generated by a program. \\
--- & The transformation & TPC publishes rules in English, not code. Third parties have written implementations. \\
\bottomrule
\end{tabular}
\end{center}

\subsection{Repositories to investigate}

\begin{center}
\small
\begin{tabular}{@{}cL{4.6cm}L{8.6cm}@{}}
\toprule
\textbf{\#} & \textbf{Repository} & \textbf{Why it matters} \\
\midrule
R1 & \href{https://github.com/abroniewski/TPC-DI-MS-SQL-Benchmark}{abroniewski/ TPC-DI-MS-SQL-Benchmark} & \textbf{Investigate this first.} Appears to contain both a source-schema script and a single SQL file with all the transformations. If true, it solves two of our three missing pieces. \\
\addlinespace[2pt]
R2 & \href{https://github.com/haydarai/tpcdi-kit}{haydarai/tpcdi-kit} & Appears to be a data generator. If it works without registering with TPC, it solves the third piece. \\
\addlinespace[2pt]
R3 & \href{https://github.com/rlsalcido24/dbtpcdi}{rlsalcido24/dbtpcdi} & A dbt implementation targeting several database systems, so its SQL should be more portable. \\
\addlinespace[2pt]
R4 & \href{https://github.com/shannon-barrow/databricks-tpc-di}{shannon-barrow/ databricks-tpc-di} & The largest and best-maintained implementation, but built for the Databricks cloud platform. May not run on a laptop. \\
\addlinespace[2pt]
R5 & \href{https://github.com/detobel36/tpc-di}{detobel36/tpc-di} & Where our target schema came from. Its transformations are in a format we cannot read easily (SSIS), so only the schema file is useful. \\
\bottomrule
\end{tabular}
\end{center}

\begin{warnbox}
The ``Why it matters'' column records expectations drawn from repository descriptions,
not verified facts --- these files have not been opened. Confirming or refuting each
expectation is the purpose of tasks A1--A3.
\end{warnbox}

% ----------------------------------------------------------------------
\subsection{TPC-DI tasks}

\begin{taskbox}
\textbf{Task A1 --- Clone the repositories and look inside} \hfill \emph{2 hours}
\end{taskbox}

\textbf{Goal.} Establish the contents of each repository.

\textbf{Steps.} From \code{\textasciitilde/benchmark}:

\begin{lstlisting}[style=sh]
cd ~/benchmark
mkdir -p tpcdi && cd tpcdi
git clone https://github.com/abroniewski/TPC-DI-MS-SQL-Benchmark.git
git clone https://github.com/haydarai/tpcdi-kit.git
git clone https://github.com/rlsalcido24/dbtpcdi.git
\end{lstlisting}

Now list what is inside each. From \code{\textasciitilde/benchmark/tpcdi}:

\begin{lstlisting}[style=sh]
find TPC-DI-MS-SQL-Benchmark -type f | head -50
find tpcdi-kit -type f | head -50
find dbtpcdi -type f -name "*.sql" | head -50
\end{lstlisting}

\textbf{Deliverable.} In \code{FINDINGS.md}, paste the file listings. Then answer:

\begin{enumerate}
\item Does \code{TPC-DI-MS-SQL-Benchmark} contain a file named
  \code{CreateDBTableSchema.sql}? What is its full path?
\item Does it contain \code{historical\_load.sql}? What is its full path?
\item How many lines does each have? Use \code{wc -l \textless filename\textgreater}.
\item Does \code{tpcdi-kit} contain a README explaining how to generate data? Does the
  README mention needing to register or download anything from tpc.org?
\end{enumerate}

\vspace{4pt}
\begin{taskbox}
\textbf{Task A2 --- Read the source schema file} \hfill \emph{2 hours}
\end{taskbox}

\textbf{Goal.} Determine whether \code{CreateDBTableSchema.sql} supplies $S$.

\textbf{Steps.} Open the file in an editor.

For each \code{CREATE TABLE} statement you find, record: the table name, the number of
columns, and the first three column names with their types.

\textbf{Example of what a good answer looks like:}

\begin{lstlisting}[style=sql]
CREATE TABLE source.Customer (
    CustomerID   INT,
    TaxID        CHAR(20),
    Status       CHAR(10),
    ...
);
\end{lstlisting}
\begin{center}\small
\begin{tabular}{@{}lcl@{}}
\toprule
Table & Columns & First three columns \\
\midrule
\code{source.Customer} & 23 & \code{CustomerID INT}, \code{TaxID CHAR(20)}, \code{Status CHAR(10)} \\
\bottomrule
\end{tabular}
\end{center}

\textbf{Deliverable.} A table in \code{FINDINGS.md} listing every source table, and a
one-line verdict: \emph{``This file does / does not give us $S$.''}

\vspace{4pt}
\begin{taskbox}
\textbf{Task A3 --- Classify the transformations} \hfill \emph{1 day. This is the most important task.}
\end{taskbox}

\textbf{Goal.} Characterise the transformation classes TPC-DI requires. This determines
whether we adopt the benchmark.

\textbf{Rationale.} If most target columns are plain copies, the benchmark will not
discriminate between methods. The distribution across transformation classes is the
deciding measure.

\textbf{The categories.} Every target column falls into exactly one of these:

\begin{center}\small
\begin{tabular}{@{}lL{7.4cm}L{5.0cm}@{}}
\toprule
\textbf{Category} & \textbf{What it means} & \textbf{Example} \\
\midrule
copy & Value taken unchanged from one source column & \code{CustomerID} \\
rename & Same value, different column name & \code{cust\_id} $\rightarrow$ \code{CustomerID} \\
retype & Value converted to a different type & text \code{'2024-01-05'} $\rightarrow$ a date \\
concat & Two or more source columns glued together & \code{first\_name || ' ' || last\_name} \\
conditional & Value depends on a test & \code{CASE WHEN rating < 'BBB' THEN 1 ELSE 0 END} \\
arithmetic & Value computed by calculation & \code{revenue / earnings} \\
lookup join & Value fetched from another table & status code \code{'ACTV'} $\rightarrow$ \code{'Active'} \\
aggregate & Value summarises many rows & \code{SUM}, \code{COUNT}, \code{MAX} \\
window & Value depends on a range of other rows & highest price in the last 52 weeks \\
key generation & A number invented by the transformation & \code{ROW\_NUMBER()} \\
SCD history & The \code{IsCurrent} / \code{EffectiveDate} / \code{EndDate} columns & --- \\
constant & A fixed value, not from the source & \code{BatchID = 1} \\
\bottomrule
\end{tabular}
\end{center}

\textbf{Steps.}

\begin{enumerate}
\item Open \code{historical\_load.sql} in VS Code.
\item Find the part that fills \code{DimCustomer}. Search for
  \code{INSERT INTO DimCustomer} or \code{DimCustomer AS}.
\item For \emph{every column} it produces, decide which category it belongs to.
\item Do the same for \code{DimAccount} and \code{FactMarketHistory}.
\end{enumerate}

\textbf{Deliverable.} A spreadsheet, \code{tpcdi\_transformations.csv}, with these columns:

\begin{center}\small
\begin{tabular}{@{}lllll@{}}
\toprule
\code{target\_table} & \code{target\_column} & \code{category} & \code{source\_columns} & \code{sql\_snippet} \\
\midrule
DimCustomer & CustomerID & copy & Customer.CustomerID & \code{c.CustomerID} \\
DimCustomer & Status & lookup join & Customer.Status & \code{JOIN StatusType st ON ...} \\
DimCustomer & IsCurrent & SCD history & --- & \code{CASE WHEN ... THEN 1} \\
DimCustomer & MarketingNameplate & conditional & NetWorth, Income, \ldots & \code{CASE WHEN ...} \\
\bottomrule
\end{tabular}
\end{center}

Plus a count of how many columns fell into each category.

\begin{notebox}
\textbf{Cross-check.} Analysis of the target schema alone gave approximately 60\% plain
copies and 40\% more complex. A materially different count from the actual SQL means one of
the two analyses is wrong; report the discrepancy.
\end{notebox}

\vspace{4pt}
\begin{taskbox}
\textbf{Task A4 --- Find out how to get the data} \hfill \emph{half a day}
\end{taskbox}

\textbf{Goal.} Determine whether we can generate $D$ without registering with TPC.

\textbf{Context.} The official generator, \code{DIGen.jar}, requires registration and a
licence agreement with TPC. Several third-party replacements exist; the question is whether
any of them works standalone.

\textbf{Steps.}

\begin{enumerate}
\item Open the README of \code{tpcdi-kit} (R2). Does it include the generator itself, or
  does it tell you to download \code{DIGen.jar} from TPC first?
\item If it includes a generator, follow its instructions and try to produce data at the
  smallest scale factor.
\item Report exactly what files were produced, their names, and their sizes. Use:
\begin{lstlisting}[style=sh]
ls -lh <output-folder>
\end{lstlisting}
\item Open one of the produced files with \code{head -5 \textless filename\textgreater}
  (shows the first 5 lines) and paste those lines into \code{FINDINGS.md}.
\end{enumerate}

If registration proves unavoidable, stop and report --- licence agreements require a
supervisor decision.

\vspace{4pt}
\begin{taskbox}
\textbf{Task A5 --- Load the target schema} \hfill \emph{1 hour}
\end{taskbox}

\textbf{Goal.} Confirm the extracted target schema loads in your environment.

\textbf{Steps.} Copy \code{tpcdi\_target\_schema.sql} into
\code{\textasciitilde/benchmark/tpcdi/}. Then create a script file. Type this command
exactly --- it writes a small Python program into a file called \code{load\_target.py}:

\begin{lstlisting}[style=sh]
cd ~/benchmark/tpcdi
pip install duckdb

cat > load_target.py << 'ENDOFFILE'
import duckdb, re

sql = open('tpcdi_target_schema.sql').read()
sql = re.sub(r'^\s*--.*', '', sql, flags=re.M)     # remove comment lines

con = duckdb.connect('tpcdi.duckdb')
for statement in sql.split(';'):
    statement = statement.strip()
    if statement:
        con.execute(statement)

count = con.execute('SELECT count(*) FROM information_schema.tables').fetchone()[0]
print('tables created:', count)
ENDOFFILE

python3 load_target.py
\end{lstlisting}



\textbf{Expected output:} \code{tables created: 20}

If you get a different number, or an error, paste it verbatim into \code{FINDINGS.md}.

% ======================================================================
\clearpage
\section{Benchmark 2: MIMIC-IV $\rightarrow$ OMOP}
\label{sec:mimic}

\subsection{What this is}

\textbf{MIMIC} is a database of real hospital intensive-care records from a Boston
hospital, released for research. \textbf{OMOP CDM} (``Common Data Model'') is a standard
structure that medical researchers have agreed on, so that studies can run the same way on
data from different hospitals.

Converting MIMIC into OMOP is a genuine, real-world migration. Somebody has already written
the conversion, and published it. That makes it a promising example for us.

\subsection{Two versions --- use the right one}

\begin{center}
\begin{tabular}{@{}L{3.6cm}L{5.4cm}L{5.4cm}@{}}
\toprule
& \textbf{\href{https://github.com/MIT-LCP/mimic-omop}{MIT-LCP/mimic-omop}} & \textbf{\href{https://github.com/MIT-LCP/mimic-iv-demo-omop}{MIT-LCP/mimic-iv-demo-omop}} \\
\midrule
Data version & MIMIC-III & MIMIC-IV, demo subset \\
Access & Requires credentialing: an online ethics course plus an application. Takes weeks. & \textbf{Open. No application needed.} \\
Language & PostgreSQL & SQL, run by a Python script \\
\midrule
\textbf{Use this one?} & Not for now & \textbf{Yes --- start here} \\
\bottomrule
\end{tabular}
\end{center}

\begin{warnbox}
MIMIC-III access requires credentialing and a supervisor decision. The demo subset is
open and sufficient to assess the benchmark's suitability.
\end{warnbox}

\subsection{What we already know (from reading documentation only)}

\begin{center}
\begin{tabular}{@{}L{4.0cm}L{11.5cm}@{}}
\toprule
\textbf{Finding} & \textbf{Detail} \\
\midrule
It is staged & The conversion happens in named steps. Tables are prefixed by role:
\code{src\_} (a snapshot of the source), \code{lk\_..\_clean} (cleaned), \code{lk\_..\_concept}
(codes mapped), \code{lk\_..\_mapped} (joined together), \code{cdm\_} (the final target).
This is useful --- each step is separately inspectable. \\
\addlinespace[2pt]
It is many-to-one & In the MIMIC-III version, one target table is often built from several
source tables. For example the \code{MEASUREMENT} table draws on five different source
tables. This is the kind of complexity we want. \\
\addlinespace[2pt]
It uses a vocabulary & Medical codes in the source (e.g.\ \code{'44054006'}) must be
translated into OMOP's numeric concept IDs (e.g.\ \code{201826}) by looking them up in a
large reference table downloaded separately. \\
\bottomrule
\end{tabular}
\end{center}

\begin{notebox}
\textbf{About the vocabulary.} At first this looked like a problem, because the answer is
not derivable from the source and target schemas alone. But if we \emph{load the vocabulary
table as part of the source database}, then looking up a code becomes an ordinary join
against a source table --- which is exactly the kind of transformation we are studying. So
this is acceptable. Just be aware the vocabulary file is large.
\end{notebox}

% ----------------------------------------------------------------------
\subsection{MIMIC tasks}

\begin{taskbox}
\textbf{Task B1 --- Clone and survey} \hfill \emph{2 hours}
\end{taskbox}

\begin{lstlisting}[style=sh]
cd ~/benchmark
mkdir -p mimic && cd mimic
git clone https://github.com/MIT-LCP/mimic-iv-demo-omop.git
cd mimic-iv-demo-omop
find . -type f -name "*.sql" | sort
find . -type f -name "*.md" | sort
\end{lstlisting}

\textbf{Deliverable.} In \code{FINDINGS.md}:

\begin{enumerate}
\item The complete list of \code{.sql} files, with their folder paths.
\item How many files are in each folder.
\item The contents of the main README, summarised in your own words in 5--10 lines:
  what does this repository claim to do, and what does it say you need to run it?
\end{enumerate}

\vspace{4pt}
\begin{taskbox}
\textbf{Task B2 --- Find out where the data comes from} \hfill \emph{half a day}
\end{taskbox}

\textbf{Goal.} Locate $D$.

\textbf{Context.} The repository holds the conversion code; the data is hosted separately
on PhysioNet. The demo subset is openly downloadable.

\textbf{Steps.}

\begin{enumerate}
\item Read the repository README and find where it says the data comes from.
\item Go to \url{https://physionet.org/content/mimic-iv-demo/} and check: can the files be
  downloaded without logging in? Look for a ``Files'' section at the bottom of the page.
\item If yes, download them into \code{\textasciitilde/benchmark/mimic/data/}.
\item List what you got and how big each file is: \code{ls -lh}.
\item Show the first 3 lines of two of the CSV files using \code{head -3}.
\end{enumerate}

\textbf{Deliverable.} A table in \code{FINDINGS.md} listing every data file, its size, its
number of columns, and its column names.

\begin{notebox}
\textbf{Also establish} whether the conversion requires the OMOP vocabulary files from
\url{https://athena.ohdsi.org/}, which need a free account. Report the requirement and the
approximate size; the account decision sits with the supervisor.
\end{notebox}

\vspace{4pt}
\begin{taskbox}
\textbf{Task B3 --- Identify $S$, $T$, and the transformation} \hfill \emph{1 day}
\end{taskbox}

\textbf{Goal.} Map the files in the repository onto our four pieces.

\textbf{Steps.} Go through the \code{.sql} files you listed in B1 and sort each one into
a bucket:

\begin{center}\small
\begin{tabular}{@{}lL{7.0cm}L{5.4cm}@{}}
\toprule
\textbf{Bucket} & \textbf{How to recognise it} & \textbf{Corresponds to} \\
\midrule
source structure & \code{CREATE TABLE} statements for MIMIC tables (names like \code{patients}, \code{admissions}) & $S$ \\
target structure & \code{CREATE TABLE} statements for OMOP tables (names like \code{person}, \code{visit\_occurrence}) & $T$ \\
transformation & \code{INSERT INTO ... SELECT ...} statements & the reference program \\
helper & Anything else --- config, cleanup, tests & --- \\
\bottomrule
\end{tabular}
\end{center}

\textbf{Deliverable.} A table in \code{FINDINGS.md}: file path, bucket, and a one-line
description. Then state clearly which of the four pieces we now have and which we are
still missing.

\vspace{4pt}
\begin{taskbox}
\textbf{Task B4 --- Classify the transformations} \hfill \emph{1 day. Most important MIMIC task.}
\end{taskbox}

\textbf{Goal.} Apply the Task A3 classification to this benchmark.

\textbf{Steps.} Pick three target tables. Suggested: \code{person} (simplest),
\code{visit\_occurrence} (medium), \code{measurement} (most complex). Find the SQL that
fills each, and classify every output column using the \emph{same category list as
Task~A3}.

Add one extra category specific to this benchmark:

\begin{center}\small
\begin{tabular}{@{}lL{11.8cm}@{}}
\toprule
vocabulary lookup & The value comes from joining the large OMOP \code{concept} reference
table to translate a medical code into a numeric ID. \\
\bottomrule
\end{tabular}
\end{center}

\textbf{Deliverable.} \code{mimic\_transformations.csv}, same columns as
\code{tpcdi\_transformations.csv}, plus counts per category.

Also answer: \textbf{for each of the three tables, how many different source tables does it
read from?} This tells us how ``many-to-one'' the conversion is, which is a property we
particularly want.

\vspace{4pt}
\begin{taskbox}
\textbf{Task B5 --- Try to run it} \hfill \emph{1--2 days. Attempt only after B1--B4.}
\end{taskbox}

\textbf{Goal.} Get the conversion running end to end, so we have a working $D$ and the
converted output side by side.

\textbf{Steps.} Follow the repository's own instructions. The code targets a specific
environment, so adaptation will be needed.

\textbf{On failure}, record in \code{FINDINGS.md}:

\begin{enumerate}
\item The exact command you ran.
\item The complete error message, copied and pasted, not summarised.
\item What you tried in response.
\end{enumerate}

\begin{notebox}
\textbf{Optional.} The classification in B4 is what decides whether we adopt this
benchmark. Cap this task at two days; if it is still blocked, report the blocker and move
on.
\end{notebox}

% ======================================================================
\clearpage
\section{What to report back}

\subsection{Deliverables checklist}

\begin{center}
\begin{tabular}{@{}clL{8.0cm}c@{}}
\toprule
& \textbf{File} & \textbf{Contents} & \textbf{From task} \\
\midrule
1 & \code{FINDINGS.md} & Everything you ran and what happened, task by task & all \\
2 & \code{tpcdi\_transformations.csv} & Every target column classified & A3 \\
3 & \code{mimic\_transformations.csv} & Every target column classified & B4 \\
4 & \code{tpcdi\_source\_schema.sql} & The source schema, if R1 provides it & A2 \\
5 & \code{sample\_data/} & A few small sample data files from each benchmark & A4, B2 \\
\bottomrule
\end{tabular}
\end{center}

\subsection{Questions the report must answer}

For \textbf{each} of the two benchmarks, answer these explicitly. Answer each directly, with the evidence that supports it.

\begin{enumerate}
\item \textbf{Do we have all four pieces?} For $S$, $D$, $T$ and the transformation:
  yes or no, and where each one is.
\item \textbf{How varied are the transformations?} Give the count per category. State the
  percentage that are plain copies.
\item \textbf{How many-to-one is it?} For your three sampled target tables, how many
  source tables does each one read?
\item \textbf{What blocks us?} Registration, licences, missing files, code that will not
  run --- list each blocker and how long it would take to clear.
\end{enumerate}

\subsection{Interpretation of results}

\begin{center}
\begin{tabular}{@{}L{5.6cm}L{9.6cm}@{}}
\toprule
\textbf{Outcome} & \textbf{What it means} \\
\midrule
Under 30\% of columns are plain copies, and tables read from 2+ sources & Excellent. Use this benchmark. \\
30--70\% plain copies & Usable. We will select the interesting subset of tables. \\
Over 70\% plain copies & Too easy on its own. We would need to combine it with another source. \\
Blocked by registration or licensing & Report it. The supervisor decides. \\
\bottomrule
\end{tabular}
\end{center}

\subsection{Working notes}

\begin{itemize}
\item \textbf{Registration and licensing decisions sit with the supervisor.} This covers
  tpc.org, PhysioNet credentialing, and Athena. Report what is required rather than
  signing up.
\item \textbf{Error messages go into the report verbatim}, not paraphrased.
\item \textbf{Negative findings are results.} ``This repository does not contain the
  transformation SQL'' resolves an open question and should be reported as such.
\item \textbf{Every claim in this document about repository contents is an expectation
  drawn from README text, not a verified fact.} Some will be wrong. Establishing which is
  the point of tasks A1--A3 and B1--B4.
\item \textbf{Flag blockers early.} A day lost to a setup problem is worth a message.
\end{itemize}

\subsection{Suggested order and timing}

\begin{center}
\begin{tabular}{@{}clL{7.6cm}c@{}}
\toprule
\textbf{Day} & \textbf{Task} & \textbf{What} & \textbf{Hours} \\
\midrule
1 & Setup & Install WSL, Ubuntu, Python, git, VS Code & 3 \\
1 & A1 & Clone TPC-DI repos, list contents & 2 \\
1 & B1 & Clone MIMIC repo, list contents & 2 \\
2 & A2 & Read the TPC-DI source schema file & 2 \\
2 & A5 & Load the target schema into DuckDB & 1 \\
2 & B2 & Find and download the MIMIC demo data & 4 \\
3 & A3 & \textbf{Classify TPC-DI transformations} & 7 \\
4 & B3 & Sort MIMIC files into buckets & 4 \\
4 & A4 & Investigate the TPC-DI data generator & 3 \\
5 & B4 & \textbf{Classify MIMIC transformations} & 7 \\
6--7 & B5 & Try to run the MIMIC conversion (optional) & 10 \\
7 & --- & Write up and present findings & 3 \\
\bottomrule
\end{tabular}
\end{center}

A3 and B4 are the deciding tasks; the remainder is preparation for them.

\end{document}
